% GENERATED FILE. Edit canonical lessons, then run npm run generate:books.
% canonical-source-hash: fnv1a64:06273003bfe95d72
% canonical-lessons: TA-C17-nanpakal-nalliravu, TA-C17-transparent-middle-synonyms

\chapter{Noon and Midnight}
\label{ch:ta-noon-and-midnight}

This chapter is generated from the canonical micro-lessons used by Language Ladder.
Each section stays independently resumable and preserves the authored prerequisite order.

\section[\ta{நண்பகல்} \ta{நள்ளிரவு}]{\ta{நண்பகல்}, \ta{நள்ளிரவு} — one old "middle" root, two survival stories}

\label{lesson:TA-C17-nanpakal-nalliravu}



\begin{quote}
\textbf{Warm-up.} {[}PAUSE 2s{]} You've just watched Kannada and Telugu each reach for Sanskrit's \emph{madhya} to name noon. Tamil doesn't reach for Sanskrit at all here — it reaches into its own, much older layer of vocabulary instead.
\end{quote}

\begin{cousinweb}
Both of today's words trace back to the same root: \textbf{\ta{நள்}} (\emph{naḷ}), a \textbf{Proto-Dravidian} word for "\textbf{middle, centre}" — a genuine synonym of the everyday \textbf{\ta{நடு}} (\emph{naṭu}, "middle") you already know, just from an older layer of the vocabulary. \ta{நள்} didn't disappear from Tamil; it survives today mainly in \textbf{literary/poetic} compounds — \textbf{\ta{நள்ளிருள்}} (\emph{naḷḷiruḷ}, "thick darkness") and \textbf{\ta{நள்ளென்} \ta{கங்குல்}} (\emph{naḷḷeṉ kaṅkul}, "deep night") — rather than in everyday colloquial speech.
\end{cousinweb}

\begin{cousinweb}
\textbf{\ta{நள்ளிரவு}} (\emph{naḷḷiravu}) = "\textbf{midnight}" = \textbf{\ta{நள்}} (\emph{naḷ}, "middle") + \textbf{\ta{இரவு}} (\emph{iravu}, "\textbf{night}," a word every Tamil speaker still uses on its own) — \ta{நள்} shows up here exactly as you'd expect. \textbf{\ta{நண்பகல்}} (\emph{naṇpakal}) = "\textbf{noon, midday}" traces to the \textbf{same} \ta{நள்} root, fused with \textbf{\ta{பகல்}} (\emph{pakal}, "\textbf{daytime}") — but here it surfaces as \textbf{\ta{நண்}-} rather than \ta{நள்}-, a sandhi-driven surface change before the following \ta{ப}. Same root, two different faces, depending on what it's fused to.
\end{cousinweb}

\subsection*{Guided Practice}
{[}PAUSE 1s{]}

\begin{itemize}
\raggedright
  \item {[}YOU SAY: "naḷ" — an old word for "middle" — then "naḷḷiravu," midnight{]}
  \item {[}YOU SAY: "naṇpakal" — noon, the same naḷ root, resurfacing as naṇ-{]}
\end{itemize}

\begin{tcolorbox}[breakable,colback=teal!4,colframe=teal!35!black,title={Wrap-up recall}]
{[}PAUSE 3s{]} What old Proto-Dravidian root do \ta{நண்பகல்} and \ta{நள்ளிரவு} both trace back to, and what does it mean? (\textbf{\ta{நள்}} \emph{naḷ}, "middle" — surviving mainly in literary compounds like \emph{naḷḷiruḷ}, "thick darkness.") Why does it surface as \textbf{naṇ-} in \emph{naṇpakal}? (A sandhi change before following \textbf{p}.) Does Tamil reach for Sanskrit here? (\textbf{No}.) Next: meet the transparent \emph{naṭu} synonyms Tamil kept beside these older compounds.
\end{tcolorbox}
\section[\ta{நடுப்பகல்}, \ta{நடு} \ta{இரவு}]{\ta{நடுப்பகல்} and \ta{நடு} \ta{இரவு} — the transparent twins}

\label{lesson:TA-C17-transparent-middle-synonyms}



\begin{quote}
\textbf{Warm-up.} {[}PAUSE 2s{]} \emph{Naṇpakal} and \emph{naḷḷiravu} preserve an older middle-root in fused forms. Tamil also kept versions whose pieces remain obvious today.
\end{quote}

\subsection*{You'll want to know: The living synonyms}
\begin{itemize}
\raggedright
  \item \textbf{\ta{நடுப்பகல்}} (\emph{naṭuppakal}) — "middle-daytime," a dictionary synonym of \emph{naṇpakal}
  \item \textbf{\ta{நடு} \ta{இரவு}} (\emph{naṭu iravu}) — "middle night," the transparent synonym of \emph{naḷḷiravu}
\end{itemize}

Both use \textbf{\ta{நடு}} (\emph{naṭu}, "middle"), the everyday word still heard in \emph{naṭuvē}, "in the middle."

Latin daughter languages often inherit one fate for an old compound: eroded or transparent. Tamil kept \textbf{both fates inside one language}—old fused forms and fresh transparent synonyms side by side.

\subsection*{Guided Practice}
{[}PAUSE 1s{]}

\begin{itemize}
\raggedright
  \item {[}YOU SAY: "naṇpakal · naṭuppakal" — two noon forms{]}
  \item {[}YOU SAY: "naḷḷiravu · naṭu iravu" — two midnight forms{]}
  \item {[}YOU POINT: \emph{naṭu} as the visible everyday "middle"{]}
\end{itemize}

\begin{tcolorbox}[breakable,colback=teal!4,colframe=teal!35!black,title={Wrap-up recall}]
{[}PAUSE 3s{]} Which everyday word builds the transparent forms? (\textbf{\ta{நடு}}, \emph{naṭu}, "middle.") What did Tamil keep simultaneously? (\textbf{Older fused compounds and transparent synonyms}.) Are any pieces Sanskrit loans? (\textbf{No}.)
\end{tcolorbox}
